% FRAGMENT A11 (drafter A: class model and main results)
% REPLACES: main.tex lines 873–884 ("\subsection{5.2 Acquisition values}" to "… records the numerical examples.")
% CHANGE: Section 5.2 becomes the nondegenerate case of the class model. Title changes from
%   "Acquisition values" to "Value classes"; the label stays.
% LABELS KEPT: sec-acquisition-values.  LABELS REMOVED/ADDED: none.
% NUMBERS: all printed numbers already appear in the manuscript (moderate vector, Table 3
%   Panel B: 0.250000 and 0.526805; value classes at p=0.5, Table 4 Panel B: 0.250000 and
%   0.522760). None is new. The moderate vector's old low cost 0.3 lay below its floor bound
%   (window [0, 0.4965)), so its entry is the same at cost zero. NUM PLACEHOLDERS: none.
% CHECKS (print-only script audit/revision_v2_2026_10_21/scripts/A_reading_n_checks.py,
%   section 4): for the value classes at the benchmark vector, (A2) margins 1.195660 and
%   0.217192; (A3) 0.017014 < 0.02 < 0.022416; part (iii) margin 0.002658 at r_2=3.6;
%   W_H=1. The displayed spread formula matches direct integration at r = 1.2, 3, 3.6.
% SOURCES: Q1_N_main.md Claims 2 (N6), 11, 12(c); SYNTHESIS.md B.1.
\subsection{5.2 Value classes}\label{sec-acquisition-values}

The reversal needs no point masses, no large class gap, and no
special prior. Take any class pair under the restrictions of
\hyperref[sec-values-preparation-and-the-sale]{Section 2.2} that
differs above the reserve. Take any prior, noise scale, and low-cost
probability. \hyperref[result-proposition-2]{Proposition 2} then
holds, with its ownership condition, on a nonempty open set of
incumbent strengths and costs. The construction puts both strengths
just above the lowest value where the classes start to differ above
the reserve. The spread vanishes near that value, so informed trading
against the weaker incumbent cannot cover its cost. The profit gap
there stays near the full class value gap. A high cost that only
favorable prices justify fits in that gap. In the point-mass case
that value is \(\ell\), so the reversal holds for every \(h>\ell\). With
\(h=2\) and \(\ell=1\), the moderate-value specification satisfies all
strict inequalities of \hyperref[result-proposition-2]{Proposition 2}.
Entry rises from 0.250000 to 0.526805.
\hyperref[paper-appendix]{Appendix A} provides the construction; Online
Appendix C.4 records the numerical examples.

Dispersed values within each class give a nondegenerate example. In the
value classes of \hyperref[sec-design]{Section 6.3}, values are uniform
with half-width \(\varepsilon_V=0.05\) around \(\ell\) and \(h\), the
prior is one half, and the investor observes only the class. At the
reserve \(p=0.5\) both classes lie above the reserve. At both benchmark
strengths the low class lies inside the incumbent's support and the high
class lies above it, so

\[
\Delta_T(r)=\frac{(r-\ell)^2+\varepsilon_V^2/3}{2r}.
\]

Dispersion within the low class adds to the spread. The benchmark vector
satisfies every strict inequality of
\hyperref[result-proposition-2]{Proposition 2} for these classes, and
preparation rises from 0.250000 to 0.522760 (Table~\ref{tab:4}, Panel
B). Equal admissibility matters. At a reserve that excludes the low
class, a stronger incumbent can narrow the spread, and
\hyperref[sec-design]{Section 6.3} reports deterrence at such a reserve.
